% from ega0-22-fr.tex
\oldpage[0\textsubscript{IV-1}]{205}

  \item[\rm(ii)] Le corollaire (22.5.9) amène à considérer la question
  suivante~: si $A$ est une $k$-algèbre noethérienne, à quelles conditions
  l'ensemble des idéaux premiers $\mathfrak p\in\operatorname{Spec}(A)$ tels que
  $A_{\mathfrak p}$ soit une $k$-algèbre formellement lisse est-il
  \emph{ouvert}~? Nous en aborderons plus tard certains cas particuliers.
\end{enumerate}
\end{remark}

\subsection{Critère jacobien de Zariski}
\label{subsection:0.22.6-fr}

\begin{theorem}[22.6.1]
\label{0.22.6.1-fr}
\textnormal{(critère jacobien de lissité formelle).}
Soient $A$, $B$ deux anneaux topologiques,
$u:A\longrightarrow B$ un homomorphisme continu faisant de $B$ une
$A$-algèbre formellement lisse, $\mathfrak J$ un idéal de $B$ (non
nécessairement fermé), $C=B/\mathfrak J$ la $A$-algèbre topologique quotient.
Pour que $C$ soit une $A$-algèbre formellement lisse, il faut et il suffit que
l'homomorphisme canonique (cf. (20.5.11.2))
\begin{equation}
\label{0.22.6.1.1-fr}
  \delta_{C/B/A}:\mathfrak J/\mathfrak J^2
  \longrightarrow\Omega^1_{B/A}\otimes_BC
  \tag{22.6.1.1}
\end{equation}
soit \emph{formellement inversible à gauche} (cf. (19.1.5)).
\end{theorem}

\begin{proof}
En effet, dire que $B$ (resp. $C$) est une $A$-algèbre formellement lisse
signifie (19.4.4) que l'on a
$\operatorname{Exalcotop}_A(B,L)=0$ (resp.
$\operatorname{Exalcotop}_A(C,L)=0$) pour tout $B$-module (resp. $C$-module)
discret $L$ annulé par un idéal ouvert de $B$ (resp. $C$). Comme par hypothèse
$\operatorname{Exalcotop}_A(B,L)=0$ pour tout $C$-module discret $L$ annulé
par un idéal ouvert de $C$, dire que $C$ est une $A$-algèbre formellement lisse
signifie donc que l'homomorphisme canonique
$\operatorname{Exalcotop}_A(C,L)\longrightarrow
\operatorname{Exalcotop}_A(B,L)$ est injectif, et le théorème résulte donc de
(20.7.8).
\end{proof}

Rappelons (19.5.3) que lorsque $B$ et $C$ sont des $A$-algèbres formellement
lisses, $\mathfrak J/\mathfrak J^2$ est un $C$-module topologique
\emph{formellement projectif}~; en outre, les homomorphismes canoniques
$\varphi_n:\mathbf S_C^n(\mathfrak J/\mathfrak J^2)
\longrightarrow\mathfrak J^n/\mathfrak J^{n+1}$ sont des
\emph{bimorphismes formels} lorsque $B$ est supposé être un anneau
\emph{préadmissible}.
\begin{corollary}[22.6.2]
\label{0.22.6.2-fr}
Les hypothèses et notations étant celles de (22.6.1), supposons de plus que
dans $B$ le carré de tout idéal ouvert soit ouvert, et que sur $\mathfrak J$ la
topologie induite par celle de $B$ soit identique à la topologie déduite de
celle de $B$ (19.0) (on notera que ces deux conditions sont vérifiées lorsque
$B$ est noethérien et sa topologie préadique ($0_{\mathrm I}$, 7.3.2), ou si la
topologie de $B$ est la topologie $\mathfrak J$-préadique). Soit
$(\mathfrak b_\lambda)$ un système fondamental d'idéaux ouverts de $B$ et
posons $B_\lambda=B/\mathfrak b_\lambda$ pour tout $\lambda$. Alors~:
\begin{enumerate}
  \item[\rm(i)] Les conditions suivantes sont équivalentes~:
  \begin{enumerate}
    \item[a)] $C$ est une $A$-algèbre formellement lisse.
    \item[b)] Pour tout $\lambda$, l'homomorphisme de $B_\lambda$-modules
    \begin{equation}
    \label{0.22.6.2.1-fr}
      (\mathfrak J/\mathfrak J^2)\otimes_BB_\lambda
      \longrightarrow\Omega^1_{B/A}\otimes_BB_\lambda
      \tag{22.6.2.1}
    \end{equation}
    déduit de $\delta_{C/B/A}$ par tensorisation, est inversible à gauche
    (autrement dit un isomorphisme sur un facteur direct de
    $\Omega^1_{B/A}\otimes_BB_\lambda$).
  \end{enumerate}
  \item[\rm(ii)] Supposons de plus que l'anneau topologique $C$ soit
  préadmissible ($0_{\mathrm I}$, 7.1.2) et

  